Krátké shrnutí: postupy pro návrh ústních zkoušek zaměřené na ověření porozumění a procesu myšlení, nikoli jen výsledku.

Praktický postup pro návrh a vedení ústní zkoušky
\begin{enumerate}
  \item Příprava otázek: krátké „strukturální“ úkoly s 2–3 doplňujícími variantami; pro každou otázku definujte očekávané kroky řešení.
  \item Požadavek na procesní artefakty: student předloží krátké zápisky/skici/OneNote pro ověření mezikroků \cite{kb_ai_101_tipu_pdf_04e4a115_page_8_1}.
  \item Průběh zkoušky: warm‑up, hlavní úloha, učitel volí 1–3 doplňující otázky zaměřené na odůvodnění a adaptace; kontrola konkrétních mezikroků.
  \item Ověření autenticity: preferujte otázky o rozhodnutích, alternativách a krátkých adaptacích; zaznamenávejte klíčové body (GDPR‑kompatibilně).
\end{enumerate}

Šablony doplňujících otázek (vyberte 1–3): „Proč jste zvolil/a tento postup? Jaké byly alternativy?“, „Jak upravit řešení, pokud se X změní na Y?“, „Který krok je nejrizikovější a proč?“

Krátká rubrika: proces a logika (0–3); reakce/adaptace (0–3); kvalita artefaktů (0–2); metakognice (0–2).

Doporučení: předem informujte studenty o pravidlech a povinných artefaktech; využijte digitální nástroje pro sběr (OneNote usnadňuje čitelnost) \cite{kb_ai_101_tipu_pdf_04e4a115_page_8_1}. Poznámka: respektujte fakultní pravidla, GDPR a interní politiky VUT.